%!TeX root=../Memory.tex

\chapter{Conclusions}
\label{cap:conclusions}

Aquest capítol tanca la memòria valorant el que s'ha entregat respecte
al que s'havia plantejat al començament del projecte. La primera part
revisa el producte final --el que s'ha fet, el que s'ha decidit no fer,
el que queda per fer més endavant i les condicions de posada en
producció i manteniment. La segona part és una reflexió de l'equip
sobre el procés: què ha funcionat, què no, com s'ha treballat
conjuntament i quin valor real aporta el resultat a les diferents
parts interessades.

%-----------------------------------------------------------------------
\section{Producte final entregat}
\label{sec:producte-final}

\subsection{Què s'ha fet}

El projecte entrega un \acl{MVP} (\ac{MVP}) \textbf{funcional i
demostrable} d'un sistema de votació descentralitzada amb les peces
següents:

\begin{itemize}
    \item \textbf{Smart contract d'Algorand} escrit en \emph{algopy}
    \cite{algorandfoundation2024puya}, amb gestió d'organitzacions
    (creació, cens d'adreces autoritzades amb càrrega massiva fins a
    set adreces per transacció), propostes amb fase d'aprovació per
    quòrum de dos terços, votació preferencial amb \emph{rànquings}
    complets, i prevenció de doble vot.

    \item \textbf{Smart contract d'Ethereum} (\texttt{NotaryContract})
    en Solidity 0.8.20, desplegable a la \emph{testnet} Sepolia, que
    implementa un protocol de consens K-de-N entre nodes
    universitaris. La llista de nodes autoritzats i el llindar K es
    congelen en obrir una elecció, garantint que cap canvi posterior
    pugui alterar els paràmetres de validació.

    \item \textbf{Sistema d'ancoratge} en Python que llegeix l'estat
    final d'una elecció des d'Algorand, produeix un hash determinista
    (JSON canònic + SHA-256) i el publica al contracte notari
    d'Ethereum. La determinació de l'ordre dels vots per direcció
    lexicogràfica garanteix que tots els nodes generen exactament el
    mateix hash si no hi ha manipulació.

    \item \textbf{Client web} amb React 19 i TypeScript, amb interfície
    completa per a la creació d'organitzacions, càrrega de cens via
    \ac{CSV}, gestió de propostes, votació preferencial amb
    reordenació per arrossegament, recompte amb el mètode de Schulze
    \cite{schulze2011newmethod} computat localment, i panell de
    verificació pública que permet consultar el vot emès per qualsevol
    adreça.

    \item \textbf{Infraestructura local reproductible} amb Docker
    Compose: un sol \texttt{docker compose up} aixeca un \emph{LocalNet}
    d'Algorand amb dos nodes (\emph{lead} i \emph{follower}),
    l'\emph{Indexer} amb la seva base de dades PostgreSQL, el servei
    Conduit, compila i desplega el contracte, i serveix el client web a
    través de Vite.

    \item \textbf{Pràctiques d'enginyeria del programari}
    consolidades: \emph{backlog} jerarquitzat amb nou èpiques i més de
    100 \emph{issues} traçades a GitHub Projects, control de versions
    amb \emph{feature branches} i revisió creuada obligatòria,
    \emph{Conventional Commits} \cite{conventionalcommits} validats
    automàticament tant en local com en CI, dos \aclp{ADR} formals i
    diagrames C4 \cite{brown2018c4} a tots els nivells.
\end{itemize}

El sistema permet executar el cicle complet d'una elecció --crear una
organització, carregar un cens, redactar una proposta, obtenir-ne
l'aprovació per quòrum, votar preferencialment, calcular el guanyador
amb Schulze i ancorar el resultat a Ethereum-- en menys de cinc minuts
sobre l'entorn de desenvolupament local.

\subsection{Què s'ha decidit no fer (i no es farà)}

Una de les decisions arquitectòniques que ha pres l'equip mereix una
explicació explícita perquè condiciona l'evolució futura del
projecte: \textbf{els nodes de la xarxa permissionada no seran els
propis votants}. L'arquitectura final assumeix que els nodes Algorand
estan operats per entitats institucionals --universitats, entitats
civiques, col·legis professionals-- i que els votants accedeixen al
sistema com a usuaris finals d'un client que es connecta a aquests
nodes a través d'una \ac{API}.

Aquesta decisió s'ha pres pels motius següents:

\begin{itemize}
    \item Operar un node Algorand requereix recursos de maquinari,
    coneixement tècnic i un ample de banda constant. Esperar que cada
    votant operi un node és tècnicament irreal i pràcticament
    excloent per a usuaris no tècnics.

    \item El model de nodes institucionals s'alinea amb el cas d'ús
    natural del projecte: cada universitat opera un node, contribueix
    a la xarxa permissionada i participa al consens K-de-N
    d'ancoratge.

    \item La distribució de confiança ja queda garantida pel fet que
    múltiples entitats independents operen els nodes; no cal que cada
    votant en sigui un per assolir el mateix grau de
    descentralització.
\end{itemize}

Aquesta orientació diferencia clarament el projecte de plataformes com
DAVINCI \cite{vocdoni2026davinci}, on l'usuari final pot participar
activament a la xarxa de \emph{sequencers}. Optar per nodes
institucionals és una elecció conscient: aposta per la
\textbf{verificabilitat pública} sense exigir una participació activa a
la infraestructura per part del votant.

\subsection{Què no s'ha fet (però es farà més endavant)}

Les funcionalitats següents queden identificades explícitament com a
treball futur. No s'han abandonat: són part del projecte a llarg
termini i corresponen a èpiques del \emph{backlog} que es van
prioritzar a la fase de planificació però que s'han hagut de moure
fora de l'abast del \ac{MVP}.

\paragraph{Criptografia avançada del vot.} Incorporar xifratge
homomòrfic ElGamal exponencial \cite{adida2008helios} i proves
\emph{Groth16} de coneixement zero \cite{cortier2019belenios} per a la
verificació de l'elegibilitat sense revelar la identitat del votant.
Aquesta línia està documentada amb detall a les èpiques E3 i E4 del
repositori \cite{demochainrepo} i requereix una dedicació
estimada de diversos mesos addicionals per ser duta a producció amb
les garanties necessàries.

\paragraph{Cerimònia distribuïda de desxifrat amb \emph{trustees}.}
Una vegada introduït el xifratge, caldrà implementar el protocol de
generació distribuïda de claus (DKG) i de desxifrat en fases amb
proves Chaum-Pedersen per evitar que cap \emph{trustee} individual
pugui revelar els vots. Aquesta funcionalitat correspon a la
\ac{US} E4-US3 (\emph{Cerimònia de desxifrat dels trustees}), encara
oberta.

\paragraph{Integració amb un \acl{IdP} (\ac{IdP}) extern.}
L'autenticació del votant via un proveïdor d'identitat (per exemple
Keycloak, o el directori d'identitats d'una universitat) i la
derivació determinista de la clau de votació Ed25519 a partir de
les credencials del votant és l'objectiu de l'èpica E6, amb les
històries d'usuari E6-US1 i E6-US2 encara obertes al repositori.

\paragraph{Mecanismes de resistència a la coacció
(\emph{receipt-freeness}).} Impedir que el votant pugui demostrar a un
tercer què ha votat, una propietat fonamental per prevenir compra de
vots i coerció. Aquest mecanisme està implementat al
\emph{state-of-the-art} per sistemes com DAVINCI i és imprescindible
per a desplegaments amb riscos socials reals.

\subsection{Posada en producció}

La transició del \ac{MVP} actual a un sistema desplegable per una
universitat real requereix una sèrie de passos addicionals:

\begin{enumerate}
    \item \textbf{Operació d'un node Algorand permissionat per
    cadascuna de les universitats participants.} Cada entitat hauria
    d'aixecar un node Algorand connectat a la xarxa permissionada del
    consorci, amb la configuració de seguretat adequada (TLS,
    \emph{firewalling}, monitoratge, còpies de seguretat).

    \item \textbf{Servei d'ancoratge dedicat per node.} Cada node
    universitari hauria d'operar el seu propi servei d'ancoratge amb
    una clau privada d'Ethereum donada d'alta a la \emph{whitelist}
    del contracte notari. El servei hauria de córrer com a \emph{daemon}
    de sistema (\texttt{systemd}) amb reinici automàtic, monitoratge
    del cost de gas i alertes operatives.

    \item \textbf{Desplegament estable a Ethereum.} Per a un
    desplegament real, el contracte \texttt{NotaryContract} hauria
    de passar de la \emph{testnet} Sepolia a la \emph{mainnet}
    d'Ethereum. Això implica costos reals de \emph{gas}, una auditoria
    de seguretat del contracte i una gestió curosa de les claus dels
    administradors.

    \item \textbf{Construcció en producció del client web.} El client
    actual s'executa amb el servidor de desenvolupament de Vite.
    Per a producció caldria una construcció estàtica (\texttt{bun run
    build}), servida via CDN o un servidor web convencional (Nginx),
    amb HTTPS i polítiques d'origen creuat adequades.

    \item \textbf{Procediments operatius compartits.} Documentació
    operativa: com es coordinen les universitats per actualitzar el
    contracte, què es fa si un node es perd, com es porten a terme les
    incorporacions i baixes de membres del consorci.
\end{enumerate}

\subsection{Manteniment}

El manteniment a llarg termini del sistema cobreix tres àmbits
diferenciats:

\paragraph{Manteniment del codi.} Actualitzacions periòdiques de les
dependències (algopy, react, web3.py, algosdk) per incorporar
correccions de seguretat i compatibilitat. La cobertura de tests
automàtics existent (\emph{p.ex.} els cinc fitxers de test unitari
del contracte i els tests del client) ha de mantenir-se i estendre's.

\paragraph{Auditoria de seguretat del contracte.} Els \emph{smart
contracts} d'Algorand i d'Ethereum són essencialment immutables un
cop desplegats. Qualsevol incidència crítica requereix migrar a una
nova versió i fer una transició coordinada. És imprescindible
sotmetre el contracte d'Algorand a una auditoria de seguretat
externa abans d'un desplegament real, especialment quan
s'incorporin les capes criptogràfiques.

\paragraph{Monitoratge operatiu.} En producció, l'estat de la xarxa
Algorand, el servei d'ancoratge i el client web s'haurien de
monitorar amb eines estàndard (Prometheus, Grafana) i amb alertes
configurades sobre indicadors crítics: lag dels nodes, fallades en
les submissions a Ethereum, despeses anòmales de \emph{gas}.

%-----------------------------------------------------------------------
\section{Lliçons apreses}
\label{sec:llicons}

Aquesta secció recull la reflexió de l'equip sobre el procés
seguit. No pretén ser exhaustiva sinó assenyalar els punts que, en
retrospectiva, han marcat de manera més clara el desenvolupament del
projecte.

\subsection{Què ha funcionat}

\paragraph{La inversió en recerca prèvia.} Dedicar el primer mes
gairebé en exclusiva a entendre l'ecosistema Algorand, el llenguatge
\emph{algopy} i els mètodes de votació ordinal va semblar lent
mentre durava. En retrospectiva, ha estat la decisió més encertada
del projecte: el contracte que es va escriure després estava lliure
d'errors conceptuals que altrament haurien requerit reescriptures
costoses. Els dos \acp{ADR} formalitzats al final d'aquesta fase
(\acl{ADR} 000 sobre \emph{algopy} i \acl{ADR} 001 sobre la
plataforma) van fixar el rumb i van evitar discussions repetides.

\paragraph{El \emph{stack} d'enginyeria del programari.} Les
pràctiques adoptades --\emph{Conventional Commits}
\cite{conventionalcommits}, \emph{feature branches}, \emph{pull
requests} amb revisió creuada obligatòria, GitHub Actions, \emph{git
hooks} locals que auto-formaten el codi-- han creat un entorn de
treball en què la qualitat del codi s'ha mantingut alta amb molt poc
esforç conscient. En particular, la regla de protecció de
\texttt{main} que impedeix \emph{mergejar} la pròpia \ac{PR} ha
forçat de manera natural la revisió creuada entre membres de
l'equip.

\paragraph{Docker Compose com a punt de pivot.} Tenir un fitxer
\texttt{docker-compose.yml} que aixecava tot el sistema (\emph{LocalNet}
Algorand, Indexer, Conduit, base de dades, client web, desplegament del
contracte) amb una sola comanda ha estat decisiu. Cada membre de
l'equip podia provar canvis dels altres en pocs minuts, i la
incorporació de noves dependències no requeria documentació
manual: era automàtica.

\paragraph{Diagrames C4 \cite{brown2018c4} mantinguts al repositori.}
Els diagrames a tres nivells (context, contenidors, components per a
cada peça) han servit com a llenguatge comú per a les decisions
arquitectòniques. Quan un membre nou ha de tocar una àrea que no és
la seva, els diagrames li donen l'orientació mínima per saber on i
què mirar.

\subsection{Què no ha funcionat}

\paragraph{L'ambició criptogràfica inicial.} Plantejar la
criptografia avançada (ElGamal homomòrfic + ZK-SNARKs + trustees) al
mateix temps que tots els altres components del sistema va resultar
inviable. Sense un mes addicional de recerca específica sobre cada
una d'aquestes primitives, l'equip no tenia les bases per
implementar-les amb correcció. La decisió de reduir l'abast a una
votació en clar amb verificabilitat pública (vegeu
secció~\ref{sec:reduccio-abast} del capítol~\ref{cap:planificacio}) va
ser dolorosa però necessària. La lliçó és coneguda però val la pena
recordar-la: \emph{millor entregar menys, però sòlid, que més i
trencat}.

\paragraph{La sincronització de canvis al contracte i al client.}
Quan el contracte canviava algun mètode o algun camp d'una estructura
ARC-4, el client havia de regenerar el fitxer \texttt{Demochain.arc56.json}
i adaptar els \emph{mappers}. Aquest cicle, encara que automatitzable
amb les eines d'AlgoKit \cite{algorandfoundation2024algokit}, va
requerir més coordinació manual de la que s'havia anticipat.

\paragraph{La gestió del temps en paral·lel amb la resta del grau.}
Tot l'equip combinava el projecte amb altres assignatures. Els
\emph{sprints} de dues setmanes han resultat, en pràctica, més curts
del que sembla quan es resten les setmanes amb càrrega d'exàmens o
entregues d'altres matèries. Una planificació amb \emph{sprints} de
tres setmanes hauria estat probablement més realista.

\subsection{El funcionament com a equip}

L'equip ha funcionat de manera notablement cohesionada. Quatre
membres amb perfils similars (Grau d'Enginyeria Informàtica, sense
experiència prèvia en \emph{blockchain}) i una distribució natural
de responsabilitats --un per àrea-- ha permès paral·lelitzar el
treball sense generar conflictes d'autoria. Les reunions setmanals,
inicialment presencials i progressivament híbrides, han servit tant
per a la coordinació tècnica com per a la presa de decisions
col·lectiva.

Cap decisió rellevant s'ha pres unilateralment. Els canvis d'abast,
la reorganització de \emph{sprints} i la incorporació o eliminació
d'eines tècniques s'han discutit en grup i s'han documentat
posteriorment (al \emph{backlog}, als \acp{ADR} o en aquesta
memòria). Aquest patró ha estat un dels actius més clars del
projecte i s'hauria de mantenir si l'equip continua treballant en
versions posteriors del sistema.

La revisió creuada de \acsp{PR}, forçada per la configuració del
repositori, ha contribuït a una comprensió compartida del codi:
encara que cada membre lideri una àrea, tots tenen un coneixement
funcional bàsic de les altres tres. Aquesta propietat és
especialment valuosa per a la continuïtat del projecte.

\subsection{El producte: orgull i autocrítica}

Globalment, l'equip considera que el producte entregat compleix
amb suficiència els objectius plantejats al
capítol~\ref{cap:introduccio}. El sistema funciona, és reproductible,
està documentat amb rigor i les decisions arquitectòniques són
defensables. Les renúncies fetes al llarg del projecte --notablement
la criptografia avançada-- s'han pres amb consciència i s'han
documentat amb transparència, sense ocultar-les ni dissimular-les.

L'autocrítica principal és que el sistema, en la seva forma actual,
encara no està preparat per un desplegament en un context amb riscos
socials reals. Sense la capa de privacitat criptogràfica i sense una
auditoria de seguretat externa del contracte, recomanar el sistema
per a unes eleccions seriós seria irresponsable. Aquesta limitació
no és un defecte amagat sinó una propietat coneguda del \ac{MVP}: el
producte actual és, com el seu nom indica, una \emph{prova de
viabilitat}, no un sistema productiu.

%-----------------------------------------------------------------------
\section{Aportació a les parts interessades}
\label{sec:aportacio-stakeholders}

Finalment, val la pena articular què aporta el projecte a cada una de
les parts interessades identificades al començament:

\paragraph{A les universitats i entitats educatives.} Una base de
codi oberta, documentada i reproductible que es pot avaluar com a
punt de partida per a un sistema propi de votació interna
(elecions de rector, claustre, representació estudiantil). El
sistema actual no és apte per a aquest desplegament sense la capa
criptogràfica, però la infraestructura --Algorand permissionat +
ancoratge a Ethereum + client web obert-- és transferible
immediatament.

\paragraph{A entitats civiques i col·lectius professionals.} La
demostració que és tècnicament possible construir un sistema de
votació descentralitzada de baix cost (les comissions d'Algorand són
de l'ordre del miliesim de cèntim per transacció
\cite{chen2017algorand}) sense renunciar a la verificabilitat
pública. El projecte mostra una alternativa viable als sistemes
comercials privatius emprats actualment per moltes entitats.

\paragraph{A la comunitat acadèmica i de recerca en votació
electrònica.} Un cas d'estudi documentat sobre la combinació
Algorand + Ethereum + Schulze. Tot i no aportar primitives
criptogràfiques noves, el projecte documenta amb rigor les
restriccions reals de construir un sistema de votació amb
\emph{blockchain} en un horitzó temporal acotat, una informació útil
per a futurs equips que afrontin un projecte similar.

\paragraph{A l'equip mateix.} Quatre estudiants amb una experiència
formativa significativa sobre tecnologies emergents (\emph{blockchain},
\emph{smart contracts}, criptografia aplicada), pràctiques
d'enginyeria del programari (GitHub Projects, \aclp{ADR}, C4,
\emph{Conventional Commits}) i treball en equip estructurat. Les
competències adquirides són transferibles immediatament a contextos
professionals.

%-----------------------------------------------------------------------
\section{Tancament i línies de treball futures}
\label{sec:tancament}

Si l'equip continués el projecte més enllà de l'assignatura,
l'ordre lògic de prioritats seria el següent:

\begin{enumerate}
    \item Incorporar la capa de \textbf{privacitat criptogràfica} (xifratge
    homomòrfic ElGamal + ZK-SNARKs Groth16) com a primer pas
    indispensable abans de qualsevol desplegament real. Aquesta és la
    barrera més clara que separa el sistema actual d'un sistema
    productiu.

    \item Integrar la \textbf{cerimònia de \emph{trustees}} amb DKG i
    desxifrat distribuït, completant el model d'amenaces de manera
    coherent amb l'arquitectura criptogràfica.

    \item Integrar un \textbf{\acl{IdP}} extern per a l'autenticació
    dels votants i la derivació de claus determinista.

    \item Sotmetre el sistema a una \textbf{auditoria de seguretat
    externa} dels dos \emph{smart contracts}.

    \item Realitzar un \textbf{desplegament pilot} amb una entitat
    universitària real per a una elecció de baix risc (per exemple,
    l'elecció dels delegats d'una assignatura) que serveixi de banc
    de proves abans d'eleccions amb implicacions més grans.
\end{enumerate}

El sistema actual representa un punt de partida sòlid des del qual
es pot construir sense haver de recomençar. Aquesta és, en última
instància, la millor mètrica de l'èxit del projecte: no és el
sistema definitiu, però és una base sobre la qual el sistema
definitiu es pot construir amb confiança.
